% chapters/requirements/rf_config.tex
% RF-15: Configuración y mantenimiento
% ============================================================================

\item \textbf{Configuración y mantenimiento.} \label{rf:config}
El sistema expone un conjunto de parámetros de configuración global que los gestores pueden
modificar en caliente para su organización, sin reinicio del servicio. Además, el sistema
ejecuta tareas de mantenimiento programadas, como el borrado automático de exámenes
archivados con una antigüedad configurable.

\begin{functional}

    % RF-15.1 ---------------------------------------------------------------
    \item \textbf{Consulta de configuración de organización.}
    \label{rf:config:query}
    El sistema debe permitir a un gestor consultar todos los parámetros de configuración de
    su organización.

    \textbf{Entrada:} petición GET al recurso de configuración de la organización.

    \textbf{Proceso:} el sistema recupera los valores actuales de todos los parámetros
    configurables para esa organización: motor \acs{ocr} por defecto, duración máxima de
    audio (en minutos), tamaño máximo de \acs{pdf} (en MB), frecuencia de borrado
    automático de exámenes archivados (en días o meses), días de antelación para la
    notificación de borrado, idioma por defecto de la organización (para correos y
    notificaciones), tiempo de expiración de \acsp{url} temporales de \acs{pdf} (en
    minutos) y permisos por defecto de cada rol en asignaturas.

    \textbf{Salida:} respuesta con el mapa completo de parámetros y sus valores actuales.

    % RF-15.2 ---------------------------------------------------------------
    \item \textbf{Modificación de configuración de organización.}
    \label{rf:config:update}
    El sistema debe permitir a un gestor modificar los parámetros de configuración de su
    organización sin necesidad de reiniciar el servicio.

    \textbf{Entrada:} petición PATCH al recurso de configuración de la organización con los
    parámetros a modificar y sus nuevos valores.

    \textbf{Proceso:} el sistema valida los nuevos valores (rangos, tipos, existencia del
    motor \acs{ocr} indicado, etc.) y los aplica. La configuración se almacena en la base
    de datos con una capa de caché en Redis para garantizar consultas de baja latencia. Al
    modificar un parámetro, se invalida la caché correspondiente. Se registra la
    modificación en el \textit{log} de auditoría con los valores anteriores y nuevos.

    \textbf{Salida:} respuesta con la configuración actualizada.

    % RF-15.3 ---------------------------------------------------------------
    \item \textbf{Borrado automático programado de exámenes archivados.}
    \label{rf:config:auto_delete}
    El sistema debe ejecutar periódicamente una tarea de borrado de exámenes archivados que
    hayan superado la antigüedad configurada por organización.

    \textbf{Entrada:} tarea periódica programada según la frecuencia configurada.

    \textbf{Proceso:} el sistema identifica, para cada organización, los exámenes cuyo curso
    archivado tiene una antigüedad superior a la configurada. Para cada examen identificado,
    se eliminan las instancias, anotaciones, calificaciones y los ficheros asociados en
    MinIO (\acs{pdf}, audios) además de los usuarios inactivos no referenciados y sus
    notificaciones. Se preservan los registros del \textit{log} de auditoría
    (nunca se borran). La operación se ejecuta de forma asíncrona en la cola de tareas. Se
    registra la operación en el \textit{log} de auditoría con el número de exámenes y
    ficheros eliminados.

    \textbf{Salida:} exámenes y ficheros eliminados. Registro de la operación en los
    \textit{logs} de aplicación y de auditoría.

    % RF-15.4 ---------------------------------------------------------------
    \item \textbf{Notificación previa al borrado automático.}
    \label{rf:config:delete_warning}
    El sistema debe notificar a los gestores de una organización con la antelación
    configurada antes de ejecutar un borrado automático.

    \textbf{Entrada:} evaluación periódica que detecta que faltan $X$~días (configurables)
    para la ejecución del próximo borrado automático en una organización.

    \textbf{Proceso:} el sistema genera una notificación dirigida a todos los gestores
    activos de la organización, informando de la fecha prevista del borrado y el número de
    exámenes que se verán afectados. Esta notificación se envía una única vez por cada
    ejecución de borrado programada.

    \textbf{Salida:} notificación \textit{in-app} creada para los gestores y, si habilitado,
    correo electrónico espejo.

    % RF-15.5 ---------------------------------------------------------------
    \item \textbf{Modificación de permisos por defecto en asignaturas.}
    \label{rf:config:default_perms}
    El sistema debe permitir a un gestor modificar los permisos por defecto con los que se
    crean los roles de usuario en una asignatura.

    \textbf{Entrada:} petición PATCH al recurso de configuración de la organización con los
    nuevos permisos por defecto, organizados por rol (\texttt{COORDINATOR},
    \texttt{TEACHER}).

    \textbf{Proceso:} el sistema valida que los nuevos valores son permisos reconocidos
    (véase \cref{tab:perms_defaults}) y los almacena. Los cambios solo afectan a las nuevas
    \texttt{SubjectMembership} creadas a partir de ese momento; las existentes no se
    modifican retroactivamente. Se registra la modificación en el \textit{log} de auditoría.

    \textbf{Salida:} respuesta con la configuración de permisos actualizada.

\end{functional}
